\chapter{Validation and Diagnostics}

The exact solution and its substituted load are
\[
 u=\sin(\pi x)\sin(\pi y)\cos(\omega t),\qquad
 f=(2\pi^2-\omega^2)\sin(\pi x)\sin(\pi y)\cos(\omega t).
\]
The free case uses \(\omega=\sqrt{2}\pi\); the forced case uses \(\omega=\pi\).
Both satisfy the boundary and initial conditions and are smooth in space
and time. These tests exercise a single standing spatial mode, not arbitrary
wave data.

\section{What the Error Curves Measure}
At every stored diagnostic time the code evaluates absolute displacement errors
\[
 e_{L^2}(t)=\left(\int_\Omega|u(t)-u_h(t)|^2\dd x\right)^{1/2},\quad
 e_{H^1}(t)=\left(\int_\Omega|\nabla u(t)-\nabla u_h(t)|^2\dd x\right)^{1/2}.
\]
The latter is a seminorm, not the full \(H^1\) norm. Neither is divided by
the exact solution norm. \texttt{integrate\_difference} and
\texttt{compute\_global\_error} sum squared cell errors using
\texttt{QGauss(degree+2)}, i.e. \(3\times3\) points for \(Q_1\).
The exact gradient is supplied explicitly. These are quadrature approximations
to the integrals, not errors sampled only at nodes.

Time-evolution plots use time on a linear horizontal axis and absolute error
on a logarithmic vertical axis. Their dips reflect the oscillatory solution
and spatial/temporal error cancellation; they are not mesh-convergence rates.

\section{Spatial Convergence}
On this shape-regular affine square mesh, standard \(Q_1\) approximation gives
\[
 \|u-I_hu\|_{L^2}\leq Ch^2|u|_{H^2},\qquad
 |u-I_hu|_{H^1}\leq Ch|u|_{H^2}.
\]
Together with Galerkin stability and suitable regularity of the solution and
its time derivatives, this motivates spatial errors of orders two and one
for the wave problem at fixed finite time. These are expected asymptotic
orders, not exact equalities or an application of the elliptic
Cea lemma directly to a time-dependent problem.
The smooth manufactured data satisfy the regularity assumptions.
The \(Q_1\) approximation space on affine cells supports these estimates
\cite{arnold}.

We define \(h=\max_C\operatorname{diam}(C)=\sqrt{2}/N\), recorded in
\texttt{parameters.csv}; it is not the side length \(1/N\).
The constant factor does not change the experimental order
\[
 p=\frac{\log(e_h/e_{\tilde h})}{\log(h/\tilde h)}.
\]
The controlled study uses \(N=12,24,48\), \(T=1\), \(\Delta t=0.00125\).
Each mesh is also run with half this time step. Changes in the final error
norms quantify time-step sensitivity; they are not rigorous bounds on the
temporal error. This check matters because total errors combine space,
time, quadrature, and linear-solver contributions.
Log--log plots show final errors versus \(h\), with slopes two and one.
The older fixed-mesh time-step table is retained as a mixed-error experiment,
not evidence of second-order time convergence.

\section{Energy and Forcing}
In both cases the code computes, with the original unconstrained matrices,
\[
 E_h^n=\tfrac12(\vect V^n)^TM\vect V^n+
       \tfrac12(\vect U^n)^TK\vect U^n.
\]
Boundary entries of the solution vectors are zero. Continuously,
\(E'(t)=(f,u_t)\), so energy is conserved for \(f=0\).
The Newmark parameters imply
\[
 E_h^{n+1}-E_h^n
 =(\vect U^{n+1}-\vect U^n)^T
   \frac{\vect F^{n+1}+\vect F^n}{2}
\]
with exact algebraic solves. Thus discrete free energy is exactly conserved
in exact arithmetic; measured drift reflects round-off and solve tolerances.
We report \(\max_n|E_h^n-E_h^0|/E_h^0\), not just a visually flat curve.

For the manufactured forced solution,
\[
 E(t)=\frac{\pi^2}{4}\cos^2(\pi t)+\frac{\pi^2}{8}\sin^2(\pi t).
\]
It decreases then increases on \([0,1]\) as the force removes and supplies
energy. Similar initial and final values alone would miss this variation.
